\documentclass{tufte-handout}

%\geometry{showframe}% for debugging purposes -- displays the margins

\usepackage{amsmath,amsthm}

% Set up the images/graphics package
\usepackage{graphicx}
\setkeys{Gin}{width=\linewidth,totalheight=\textheight,keepaspectratio}
\graphicspath{{graphics/}}

\title{Math 550 Notes}
\author{Shan Gao}
\date{Winter 2024}  % if the \date{} command is left out, the current date will be used

% The following package makes prettier tables.  We're all about the bling!
\usepackage{booktabs}

% The units package provides nice, non-stacked fractions and better spacing
% for units.
\usepackage{units}

% The fancyvrb package lets us customize the formatting of verbatim
% environments.  We use a slightly smaller font.
\usepackage{fancyvrb}
\fvset{fontsize=\normalsize}

% Small sections of multiple columns
\usepackage{multicol}

% Provides paragraphs of dummy text
\usepackage{lipsum}

% These commands are used to pretty-print LaTeX commands
\newcommand{\doccmd}[1]{\texttt{\textbackslash#1}}% command name -- adds backslash automatically
\newcommand{\docopt}[1]{\ensuremath{\langle}\textrm{\textit{#1}}\ensuremath{\rangle}}% optional command argument
\newcommand{\docarg}[1]{\textrm{\textit{#1}}}% (required) command argument
\newenvironment{docspec}{\begin{quote}\noindent}{\end{quote}}% command specification environment
\newcommand{\docenv}[1]{\textsf{#1}}% environment name
\newcommand{\docpkg}[1]{\texttt{#1}}% package name
\newcommand{\doccls}[1]{\texttt{#1}}% document class name
\newcommand{\docclsopt}[1]{\texttt{#1}}% document class option name
%--------Theorem Environments--------
%theoremstyle{plain} --- default
\newtheorem{thm}{Theorem}
\newtheorem{cor}[thm]{Corollary}
\newtheorem{prop}[thm]{Proposition}
\newtheorem{lem}[thm]{Lemma}
\newtheorem{conj}[thm]{Conjecture}
\newtheorem{quest}[thm]{Question}

\theoremstyle{definition}
\newtheorem{defn}[thm]{Definition}
\newtheorem{defns}[thm]{Definitions}
\newtheorem{con}[thm]{Construction}
\newtheorem{exmp}[thm]{Example}
\newtheorem{exmps}[thm]{Examples}
\newtheorem{notn}[thm]{Notation}
\newtheorem{notns}[thm]{Notations}
\newtheorem{addm}[thm]{Addendum}
\newtheorem{exer}[thm]{Exercise}

\theoremstyle{remark}
\newtheorem{rem}[thm]{Remark}
\newtheorem{rems}[thm]{Remarks}
\newtheorem{warn}[thm]{Warning}
\newtheorem{sch}[thm]{Scholium}



\begin{document}

\maketitle% this prints the handout title, author, and date

\begin{abstract}
\noindent Some Math 550 notes written by me
\end{abstract}


\section{Convexity}

In these types of problems, we primarily work with finite subsets of $\mathbb{R}^d$. 
\begin{defn}[Linear subspace]
	A linear subspace of $\mathbb{R}^d$ is a collection of vectors closed under sums and multiplication by scalars. 
\end{defn}

\begin{defn}[Linear hull]
	For $S \subseteq \mathbb{R}^d$, the linear hull $\langle S \rangle$ is the smallest linear subspace containing $S$, i.e, it's the intersection of all linear subspaces containing $S$. Easy to check that it's equivalent to 
	$$ \langle S \rangle = \{ \alpha_1v_1 + \alpha_2 v_2 + \ldots + \alpha_n v_n \colon v_1,\ldots,v_n \in S, \alpha_1,\ldots,\alpha_n \in \mathbb{R}\}$$
	
\end{defn}
Noting that in finite dimensional case, $\mathbb{R}^d$, only need $d$ (at most) vectors to span the whole space. In the sense that if you have a collection of vectors whose size is greater than $d$, the vectors are automatically linearly dependent. 
\begin{defn}[Linearly dependent]
	A collection $\{v_1,v_2,\ldots,v_n\}$ of vectors are linearly dependent if some $v_i$ is a linear combination of vectors in $\{v_1,\ldots,v_n\} - \{v_i\}$ or $v_i \in \langle \{v_1,\ldots,v_n\} \rangle$. 
\end{defn}
\begin{defn}[Affine subspace]
	In $\mathbb{R}^3$, planes are not necessarily passing through the origin, 
\end{defn}

\bibliography{notes}
\bibliographystyle{plainnat}





\end{document}
